\section{What the results tell us}
Taken together, the numbers describe a detector that works in the conditions it was trained under and struggles when those conditions change. In the source domain the Autoencoder separates normal from anomalous sound with an AUC around 0.65, which is in the range one expects from a simple reconstruction-based baseline on recent DCASE data. In the target domain it is close to chance, and for several machines below it. The explanation is in the method itself. The Autoencoder is trained to rebuild normal source-domain windows and nothing else; any change in the acoustic surroundings increases its reconstruction error whether or not the machine is healthy. The ten target clips that Model~B sees are about one percent of its training data, enough to shift the standardisation statistics and the decoder a little, and that is what we observe: a small, mostly positive improvement that does not remove the problem.

The comparison of scoring methods is instructive. Residual-space Mahalanobis distance performed a little better than plain MSE because it reweights the 256 residual dimensions by how variable they are on normal data; bands that are always reconstructed poorly stop dominating the score. Latent-space distance failed because nothing in the training objective asks the bottleneck to be discriminative. Anomalous windows are reconstructed badly, but their codes are not unusual, so a distance measured in code space sees nothing.

\section{On threshold calibration}
The finding we consider most useful for anyone building on this work concerns the thresholds. The problem statement, like the official baseline, uses every normal recording for training, which leaves no independent normal recording for calibration. Fitting the threshold to training scores is then the only option, and Section~\ref{sec:rates} shows what it costs: in the target domain almost half of the healthy recordings are flagged. Threshold-free metrics such as AUC and pAUC are untouched by this, but any operating point is. The Part~1 notebook contains an optional cross-fitted calibration stage, in which fold models score recordings they never trained on, that would provide independent normal scores; we did not run it for the reported results, and it is the first thing we would do with more time.

\section{Limitations}
\begin{enumerate}
  \item \textbf{No independent validation set.} All 990 or 1\,000 normal recordings train each model. Overfitting cannot be assessed, and thresholds calibrated on training scores are optimistic (34 to 36\,\% false positives on normal test recordings instead of the intended 10\,\%).
  \item \textbf{Modest detection performance.} The best average AUC is 0.573 and target-domain AUC is near or below 0.5. The system is a baseline, not a detector one could deploy.
  \item \textbf{Not comparable with the official baseline.} Two-channel features, standardisation and the residual-space statistics mean that our numbers cannot be set against published baseline results.
  \item \textbf{Channel semantics.} The two channels are used with neutral labels. Nothing in this work establishes which physical microphone each channel represents, and a stereo file alone does not prove that two separate microphones were used.
  \item \textbf{One seed, one run.} Each model was trained once with seed 42. The variance between seeds was not measured, so differences of a few hundredths of AUC between configurations are not statistically established.
  \item \textbf{Selective scoring not evaluated.} Distances to the target-domain and all-training statistics were stored for every recording, but a selective score (the minimum over domain statistics, as in the official baseline) was not evaluated as a method of its own.
  \item \textbf{Application testing.} The interface is inference-verified but not user-tested on real recordings within the submitted evidence, and the application package omits the requirements file that its instructions mention.
  \item \textbf{Earlier prototype not retained.} An earlier ToyCar prototype trained with a 10\,\% hold-out was not available to the final pipeline run and is not included; all reported ToyCar models follow the \texttt{spec\_full} protocol.
\end{enumerate}
